\documentclass[11pt,onecolumn]{IEEEtran}
\usepackage{amsmath,amssymb,amsthm}
\usepackage{graphicx}
\usepackage[hidelinks]{hyperref}

\newtheorem{theorem}{Theorem}
\newtheorem{proposition}[theorem]{Proposition}
\newtheorem{corollary}[theorem]{Corollary}
\newtheorem{lemma}[theorem]{Lemma}
\newtheorem{remark}[theorem]{Remark}
\newtheorem{definition}[theorem]{Definition}

\newcommand{\E}{\mathbb{E}}
\newcommand{\R}{\mathbb{R}}
\newcommand{\tr}{\operatorname{tr}}
\newcommand{\Cov}{\operatorname{Cov}}
\newcommand{\hY}{\widehat{Y}}
\newcommand{\hYv}{\hat{\mathbf Y}}
\newcommand{\RW}{\mathcal{R}\!\mathcal{W}}

\begin{document}

\title{Rate and Conditional Content for\\ Gaussian Vector Sources with Encoder-Observed Context}

\author{Andrew~H.~Bond,~\IEEEmembership{Senior Member,~IEEE}%
\thanks{A. H. Bond is with the Department of Computer Engineering, San Jos\'e
State University, San Jos\'e, CA 95192 USA (e-mail: andrew.bond@sjsu.edu).
The manuscript sources and the check scripts of Appendix~\ref{app:numver}
are in the geometric-observation repository of the author.}}

\markboth{Manuscript}{Bond: Rate and Conditional Content for Gaussian Vector Sources}

\maketitle

\begin{abstract}
A stored description has a rate, paid by the decoder, and a conditional
content, the least rate at which a party holding side information can
re-encode it. For a scalar described variable a companion paper solved
the problem in closed form. This paper treats a described vector. An
encoder observes a jointly Gaussian source made of a described vector and
a correlated context, a decoder reconstructs the vector from the
description alone under a total squared-error budget, and a third party
holds a noisy copy of the context. We prove the coding theorem for this
source: the region of jointly achievable rate and conditional content is
the single-letter region, and Gaussian descriptions, indexed by their
error covariance, exhaust it. When the coordinates carry independent
contexts, the region is a union of sums of scalar regions, and the
content-optimal split of the distortion equalizes slopes that we compute
in closed form. That split is not reverse water-filling, and the
rate-optimal and content-optimal descriptions can describe disjoint sets
of coordinates. When the coordinates are coupled, the minimal conditional
content is a determinant maximization in the error covariance left to the
side-information holder. In a one-dimensional regime, certified by an
explicit test, it has a closed form in the largest generalized eigenvalue
of a matrix pencil in which the distortion budget enters, and the direction
the description reads turns with the budget unless it reads a principal
axis of the described vector that is uncorrelated with the context. Principal
components, reverse water-filling, and the Gaussian information
bottleneck all read fixed directions.
\end{abstract}

\begin{IEEEkeywords}
Conditional mutual information, Gaussian vector sources, rate--distortion
theory, side information, water-filling.
\end{IEEEkeywords}

\section{Introduction}
\label{sec:intro}

A description of data stored for one party is often kept by another. The
decoder that reconstructs from the description pays its rate. A party
that keeps the description while holding side information pays its
conditional entropy given that side information, the least rate at which
that party can re-encode it. The companion paper \cite{bond2026tit}
studies this pair of costs when an encoder observes a jointly Gaussian
pair, a scalar described variable $Y$ and a correlated context $V$, and a
third party holds a noisy copy $S$ of the context. It proves the coding
theorems, evaluates the minimal conditional content in closed form,
characterizes the region of jointly achievable rate and conditional
content, and shows that the two costs are minimized by different
descriptions.

The present paper takes the described variable to be a vector. A record
in a store is rarely one number, and the question becomes which parts of
it a description should resolve. The classical answer for the rate,
reverse water-filling \cite{berger1971}, describes the principal
components of the described vector in decreasing order of variance, along
fixed directions. We show that the answer for the conditional content is
different in kind, and we organize the results around three questions.

The first is whether the single-letter region still has operational
meaning. Theorem~\ref{thm:op} proves the coding theorem for a vector
described variable, extending the Gaussian operational theorem of
\cite{bond2026tit}, and Proposition~\ref{prop:posterior} shows that
Gaussian descriptions exhaust the region and that each is summarized by
its error covariance. The holder's side information enters that summary
as added precision.

The second is how a content-optimal description splits a distortion budget
across coordinates that carry independent contexts. Lemma~\ref{lem:slope}
computes the slope and curvature of the scalar function of
\cite{bond2026tit} in closed form and proves it strictly convex, and
Theorem~\ref{thm:parallel} shows that the content-optimal split equalizes
these slopes. It begins describing a coordinate in decreasing order of
$\sigma_j^2(1+\rho_j^2/\tau_j^2)$ rather than of the variance
$\sigma_j^2$, so a small coordinate whose context the holder reads well is
described first, and the rate-optimal and content-optimal descriptions can
describe disjoint sets of coordinates
(Corollary~\ref{cor:parallelmisalign}).

The third is what happens when the coordinates are coupled, so that no
change of coordinates separates them. Theorem~\ref{thm:maxdet} shows that
the minimal conditional content is the value of a determinant
maximization under linear matrix inequalities in the error covariance that
remains for the holder, with a unique optimum.
Theorem~\ref{thm:onedim} gives the optimum in closed form whenever the
optimal description has one dimension, with an explicit certificate for
that regime, in terms of the largest generalized eigenvalue of a matrix pencil that
holds the source covariance, the conditional covariance left after the
side information, and the distortion budget. Its eigenvector is the
direction the description reads, and Theorem~\ref{thm:turning} shows that
this direction turns as the budget changes unless it is a principal axis
of the described vector that the context does not touch. For a scalar
described variable the same pencil returns the closed form of
\cite{bond2026tit}.

\subsection*{Related work}

Independent Gaussian components and reverse water-filling are classical
\cite{berger1971}, and Section~\ref{sec:parallel} uses the same
decomposition; what differs is the allocation it produces. Parameterizing
a Gaussian description by its error covariance is a standard device in
vector Gaussian source coding \cite{rahman2012}, and
Section~\ref{sec:coupled} applies it, with the determinant maximization of
\cite{vandenberghe1998} as the resulting program. The Gaussian information
bottleneck \cite{chechik2005} projects onto fixed eigenvectors of a
normalized conditional covariance, with its tradeoff parameter setting
only how many and at what scale; the content-optimal description here
reads directions that move with the budget. The single-letter functional
is Steinberg's common-reconstruction functional \cite{steinberg2009} on
the source augmented by the context, as noted for finite alphabets by
Lapidoth, Mal\"ar, and Wigger \cite{lapidoth2014}. The closest Gaussian
counterpart, Lu and Xu's vector Gaussian Heegard--Berger problem with
common reconstructions \cite{luxu2021}, constrains the error covariance
of the whole source, context included, under a hypothesis that excludes an
unpenalized context; \cite[Appendix~C]{bond2026tit} shows that its closed
form does not continue to that limit even for scalar $Y$.

The paper is organized as follows. Section~\ref{sec:model} sets up the
model and proves the coding theorem. Section~\ref{sec:parallel} treats
independent coordinates and Section~\ref{sec:coupled} coupled ones.
Section~\ref{sec:discussion} lists what remains open, and
Appendix~\ref{app:numver} describes the numerical checks.

\section{Model and the Coding Theorem}
\label{sec:model}

All logarithms are base two unless written $\ln$. The encoder observes
i.i.d.\ copies of $T=(Y,V)$, jointly Gaussian and centered, with
$Y\in\R^m$, $V\in\R^r$, $p=m+r$, and $\Sigma_T\succ0$. The third party
holds $S=V+U$ with $U\sim\mathcal N(0,\Sigma_U)$, $\Sigma_U\succ0$,
independent of $T$ and of all channel randomness. Write $F=[\,I_m\ 0\,]$
and $H=[\,0\ I_r\,]$ for the coordinate selections, so $Y=FT$ and $V=HT$,
and
\[
J_S:=H^{\mathsf T}\Sigma_U^{-1}H,
\qquad
\Sigma_{T\mid S}=(\Sigma_T^{-1}+J_S)^{-1},
\]
the information that $S$ carries about $T$ and the conditional covariance
it leaves. A code consists of an encoder $f_n$ of $T^n$ into a finite set
$\mathcal M_n$ and a decoder $g_n:\mathcal M_n\to(\R^m)^n$. With
$M=f_n(T^n)$ and $\hYv^n=g_n(M)$, its rate, distortion, and conditional
content are
\[
R_n=\tfrac1n\log|\mathcal M_n|,\qquad
D_n=\tfrac1n\sum_{i=1}^n\E\|Y_i-\hYv_i\|^2,\qquad
L_n=\tfrac1nH(M\mid S^n),
\]
and a pair $(R,L)$ is achievable at distortion $D$ when a sequence of codes
has $\limsup R_n\le R$, $\limsup L_n\le L$, and $\limsup D_n\le D$; the
closure of the achievable set is $\RW(D)$. These are the definitions of
\cite[Section~II]{bond2026tit} with a vector reproduction. Throughout,
$0<D<\tr\Sigma_Y$; at $D\ge\tr\Sigma_Y$ the zero reproduction gives
$R=L=0$. Write $R_{\min}(D)$ and $L(D)$ for the least $R$ and the least $L$
over $\RW(D)$, the minimal rate and the minimal conditional content; by
Theorem~\ref{thm:op} they are the minima of $I(T;\hYv)$ and
$I(T;\hYv\mid S)$ over admissible channels, and a channel attaining one of
them is called rate-optimal or content-optimal.

A test channel $p(\hat{\mathbf y}\mid t)$ is admissible when
$\E\|Y-\hYv\|^2\le D$, and the Markov chain $\hYv-T-S$ holds because the
description is formed without $S$. For a channel, let $\Sigma_{e0}$ and
$\Sigma_e$ be the error covariances of the best linear estimates of $T$
from $\hYv$ and from $(\hYv,S)$. We use one result of \cite{bond2026tit}.

\begin{lemma}[{\cite[Lemma~12 and Remark~13]{bond2026tit}}]
\label{lem:gauss}
Let $\hYv\in\R^m$ be centered with finite second moments, let $A$ be its regression
matrix on $T$ and $N'=\hYv-AT$ its residual, and suppose
$\Cov(N',S)=0$. Then
$I(T;\hYv)\ge\tfrac12\log(\det\Sigma_T/\det\Sigma_{e0})$ and
$I(T;\hYv\mid S)\ge\tfrac12\log(\det\Sigma_{T\mid S}/\det\Sigma_e)$, and the
jointly Gaussian channel $AT+N$ with $N\sim\mathcal N(0,\Cov N')$
independent of $(T,S)$ has the same second moments and meets both bounds
with equality. If $\Cov N'$ is singular, either both coordinates are
infinite or a nonzero combination $\eta^{\mathsf T}\hYv$ vanishes and, after an
invertible linear map of $\hYv$, may be dropped as a coordinate.
\end{lemma}

For an admissible channel, $N'$ is uncorrelated with $T$ by construction
and with $U$ by independence, hence with $S=V+U$, so the lemma applies.

\begin{proposition}\label{prop:posterior}
The union of the quadrants above
$\bigl(I(T;\hYv),I(T;\hYv\mid S)\bigr)$ over admissible channels equals
the union, over symmetric matrices $\Sigma_{e0}$ with
$0\prec\Sigma_{e0}\preceq\Sigma_T$ and $\tr(F\Sigma_{e0}F^{\mathsf T})\le D$,
of the quadrants above
\begin{equation}
\bigl(r(\Sigma_{e0}),\ \ell(\Sigma_{e0})\bigr)
=\Bigl(\tfrac12\log\frac{\det\Sigma_T}{\det\Sigma_{e0}},\
\tfrac12\log\frac{\det\Sigma_{T\mid S}}{\det\Sigma_e}\Bigr),
\qquad
\Sigma_e^{-1}=\Sigma_{e0}^{-1}+J_S .
\label{eq:posterior}
\end{equation}
\end{proposition}

\begin{proof}
Given an admissible channel, centering $\hYv$ leaves both coordinates
unchanged and does not raise the distortion. If $\Cov N'$ is singular with
infinite coordinates the channel contributes no quadrant; otherwise
Lemma~\ref{lem:gauss}, after dropping vanishing combinations until
$\Sigma_N\succ0$, replaces it by the jointly Gaussian channel $\hYv=AT+N$
with no larger distortion and no larger coordinates. Then
$\Sigma_{e0}=\Cov(T\mid\hYv)$ has
$\Sigma_{e0}^{-1}=\Sigma_T^{-1}+A^{\mathsf T}\Sigma_N^{-1}A$, so
$0\prec\Sigma_{e0}\preceq\Sigma_T$ and $I(T;\hYv)=r(\Sigma_{e0})$. The pair
$(\hYv,S)$ is a linear observation of $T$ with independent noises $(N,U)$,
so $\Sigma_e^{-1}=\Sigma_{e0}^{-1}+J_S$ and $I(T;\hYv\mid S)=\ell(\Sigma_{e0})$.
The distortion is at least $\E\|Y-\E[Y\mid\hYv]\|^2=\tr(F\Sigma_{e0}F^{\mathsf T})$.
Conversely, given $\Sigma_{e0}$, factor
$\Sigma_{e0}^{-1}-\Sigma_T^{-1}=G^{\mathsf T}G$, let $W=GT+N$ with
$N\sim\mathcal N(0,I)$ independent of $(T,U)$, and reproduce $\hYv=\E[Y\mid W]$. Then
$\Cov(T\mid W)=\Sigma_{e0}$, the distortion is
$\tr(F\Sigma_{e0}F^{\mathsf T})\le D$, $\hYv-T-S$ holds, and since $\hYv$
is a function of $W$, data processing bounds its coordinates by those of
$W$, which are \eqref{eq:posterior}.
\end{proof}

\begin{theorem}\label{thm:op}
For the source of this section and every $0<D<\tr\Sigma_Y$, $\RW(D)$
equals the union of Proposition~\ref{prop:posterior}.
\end{theorem}

\begin{proof}
The proof of the scalar Gaussian operational theorem
\cite[Theorem~29]{bond2026tit} carries over. We give each step and what
changes.

\emph{Converse.} For any code, $\log|\mathcal M_n|\ge I(T^n;\hYv^n)\ge\sum_i
I(T_i;\hYv_i)$ and $H(M\mid S^n)\ge I(T^n;\hYv^n\mid S^n)\ge\sum_i
I(T_i;\hYv_i\mid S_i)$, by the chain rule and the independence of the
pairs $(T_i,S_i)$, exactly as in \cite[Theorem~29]{bond2026tit}; neither
step uses differential entropy or the dimension of the letters, and every
term is finite, being at most $I(T_i;S_i)+\log|\mathcal M_n|$ with $I(T_i;S_i)=\tfrac12\log\det(\Sigma_V+\Sigma_U)/\det\Sigma_U<\infty$.
The mixture of the per-letter channels is admissible at distortion $D_n$,
and by the convexity of both coordinates in the channel
\cite[Lemma~3 and Remark~4]{bond2026tit}, $(R_n,L_n)$ lies in the union at
$D_n$.

\emph{Closedness.} Let $(R_j,L_j)\to(R,L)$ lie in the union at levels
$D_j$ with $\limsup_jD_j\le D$, and pick $\Sigma_j$ feasible at $D_j$ with
$r(\Sigma_j)\le R_j$ and $\ell(\Sigma_j)\le L_j$. The $\Sigma_j$ lie in the
compact set $0\preceq\Sigma\preceq\Sigma_T$, and
$\det\Sigma_j\ge\det\Sigma_T\,2^{-2R_j}$ is bounded away from zero, so a
subsequence converges to some $\Sigma^\star\succ0$, feasible at $D$ by
continuity of the trace. The functions $r$ and $\ell$ are continuous on
positive definite matrices, so $(R,L)$ lies above
$(r(\Sigma^\star),\ell(\Sigma^\star))$. This gives closedness at fixed
distortion and the passage from $D_n$ to $D$ in the converse.

\emph{Achievability.} Fix an admissible channel with
$I(T;\hYv)<\infty$; then $\E\|\hYv\|^2\le2\tr\Sigma_Y+2D<\infty$. Steps A,
B, and C of \cite[Theorem~29]{bond2026tit} apply with these changes. In
Step A the side information is quantized coordinatewise, as there. In
Step B the source $T$ and the reproduction $\hYv$ are quantized
coordinatewise at level $k$ (we write $k$ for the quantizer index of
\cite[Theorem~29, Step~B]{bond2026tit}, and $q_\eta$, $S_\eta$ for the
side-information quantizer of Step~A, to avoid clashes with $m=\dim Y$),
so $\hYv_k=q_k(\hYv)$ takes finitely many values in $\R^m$ and
$\hYv_k\to\hYv$ in $L^2$ coordinate by coordinate; the surrogate distortion
is
$d_k(t_k,\hat{\mathbf y})=\E[\|Y-\hat{\mathbf y}\|^2\mid T_k=t_k]
=\tr\Cov(Y\mid T_k=t_k)+\|\E[Y\mid T_k=t_k]-\hat{\mathbf y}\|^2$, which is
bounded for each $k$. The identification of the true and surrogate
distortions uses memorylessness only. The convergence of the surrogate
distortion to $\E\|Y-\hYv\|^2$ uses the martingale convergence
$\E[Y\mid T_k]\to Y$ in $L^2$, applied coordinatewise. The weak
convergence of the surrogate law of $(\hYv_k,S_\eta)$ is the argument
there with test functions bounded on $\R^m\times q_\eta(\R^r)$ and
continuous in the first argument, followed by the lower semicontinuity of
mutual information \cite{posner1975}. Step C is unchanged. The
corner $\bigl(I(T;\hYv),I(T;\hYv\mid S)\bigr)$ at distortion $D$ is
therefore achievable, and with the converse and
Proposition~\ref{prop:posterior} the region is as stated.
\end{proof}

\section{Parallel Described Variables}
\label{sec:parallel}

Suppose the coordinates carry their own contexts and side information:
$r=m$, and the triples $(Y_j,V_j,S_j)$, $j=1,\dots,m$, are independent,
with $\operatorname{Var}Y_j=\sigma_j^2>0$, $\operatorname{Var}V_j=1$,
$\E[Y_jV_j]=\rho_j\sigma_j$, $|\rho_j|<1$, and $S_j=V_j+U_j$,
$U_j\sim\mathcal N(0,\tau_j^2)$, $\tau_j^2>0$. Each triple is an instance
of the scalar model of \cite{bond2026tit} up to the scale of $Y_j$. Write
$\ell_j(\delta)=\tfrac12\log g^\star(\delta)$ for the scalar function of
\cite[Theorem~14(a)]{bond2026tit} at $(\rho_j,\tau_j^2)$, where
$g^\star$ is the larger root of
\[
P(g)=\delta sg^2-(\delta+s-\rho^2)g+(1-\rho^2),\qquad s=1+\tau^2,
\]
extended by $\ell_j(\delta)=0$ for $\delta\ge1$, and
$L_j(D_j)=\ell_j(D_j/\sigma_j^2)$. Write $\RW_j(D_j)$ for the
single-letter region of coordinate $j$ at distortion $D_j$, the region
of \cite[Theorem~14(b)]{bond2026tit} at normalized distortion
$D_j/\sigma_j^2$ when $0<D_j<\sigma_j^2$, the quadrant $[0,\infty)^2$
when $D_j\ge\sigma_j^2$, and empty when $D_j=0$.

\begin{lemma}\label{lem:slope}
Fix $|\rho|<1$ and $\tau^2>0$, and for $0<\delta<1$ let $g=g^\star(\delta)$,
$k=gs-1$, and
\[
\Delta=\sqrt{(\delta+s-\rho^2)^2-4\delta s(1-\rho^2)}>0.
\]
Then
$\ell(\delta)=\tfrac12\log g^\star(\delta)$ is twice differentiable on
$(0,1)$, with
\begin{equation}
\ell'(\delta)=-\frac{k}{2\ln2\,\Delta},
\qquad
\ell''(\delta)=\frac{(1-\rho^2)k^3
+\rho^2\tau^2g^2\bigl(\tau^4g+\tau^2(2g+1)+g+1\bigr)}{2\ln2\;g\,\Delta^3}>0 .
\label{eq:slope}
\end{equation}
Hence $\ell$ is strictly convex and strictly decreasing on $(0,1)$,
$\ell'(\delta)\to-\infty$ as $\delta\to0$, and
$\ell'(\delta)\to-\tau^2/\bigl(2\ln2\,(\rho^2+\tau^2)\bigr)$ as
$\delta\to1$.
\end{lemma}

\begin{proof}
The root is simple: $\partial_gP=2\delta sg-(\delta+s-\rho^2)=\Delta$ at
the larger root. By the implicit function theorem $g^\star$ is smooth, and
differentiating $P(g^\star(\delta))=0$ with $\partial_\delta P=gk$ gives
$g'=-gk/\Delta$ and $\ell'=g'/(2\ln2\,g)=-k/(2\ln2\,\Delta)$.
Differentiating once more with $k'=sg'$ and $\Delta'=2sg-1+2\delta sg'$,
and eliminating $\delta$ through $P(g)=0$, which reads
$\delta gk=(s-\rho^2)g-(1-\rho^2)$, gives the second formula, an identity
of polynomials in $(g,s,\rho^2)$ that was also confirmed by computer
algebra (Appendix~\ref{app:numver}). The first term of its numerator is
positive and the second nonnegative, so $\ell''>0$. As $\delta\to0$,
$g^\star\to\infty$ while $\Delta\to s-\rho^2>0$, so $\ell'\to-\infty$. As
$\delta\to1$, $g^\star\to1$ since $P(1)=(\delta-1)\tau^2\to0$, so
$k\to\tau^2$ and $\Delta\to|s-(1-\rho^2)|=\rho^2+\tau^2$.
\end{proof}

At $\rho=0$ the lemma returns $\ell'(\delta)=-1/(2\ln2\,\delta)$, the
slope of the classical function. Define the \emph{activation level} of
coordinate $j$,
\begin{equation}
\kappa_j:=-L_j'(\sigma_j^2-)
=\frac{\tau_j^2}{2\ln2\;\sigma_j^2\,(\rho_j^2+\tau_j^2)},
\label{eq:kappa}
\end{equation}
which for the classical function would be $1/(2\ln2\,\sigma_j^2)$.

\begin{theorem}\label{thm:parallel}
Under the model of this section:

\emph{(a) The region.}
\[
\RW(D)=\bigcup_{D_1+\cdots+D_m\le D,\ D_j\ge0}
\bigl(\RW_1(D_1)+\cdots+\RW_m(D_m)\bigr),
\]
a union of Minkowski sums, attained by product channels that describe
each coordinate by its own scalar channel.

\emph{(b) The content endpoint.} The minimal conditional content is
$L(D)=\min\sum_jL_j(D_j)$ over $D_j\ge0$ with $\sum_jD_j\le D$, and the
minimizing allocation is unique. For $\theta>0$ let
$D_j(\theta)=\sigma_j^2$ if $\theta\le\kappa_j$, and otherwise let
$D_j(\theta)\in(0,\sigma_j^2)$ solve $-L_j'(D_j)=\theta$. The optimal
allocation is $D_j=D_j(\theta^\star)$ with $\theta^\star$ the unique level
at which $\sum_jD_j(\theta^\star)=D$, attained by the product of the
channels of \cite[Theorem~14(a)]{bond2026tit} for the described
coordinates, rescaled by $\sigma_j$, and $\hY_j=0$ for the others.

\emph{(c) The order of description.} As $D$ decreases from
$\sum_j\sigma_j^2$, the content-optimal description begins to describe
coordinate $j$ when $\theta^\star$ exceeds $\kappa_j$, that is, in
decreasing order of $\sigma_j^2(1+\rho_j^2/\tau_j^2)$, while the
rate-optimal description proceeds in decreasing order of $\sigma_j^2$.
\end{theorem}

\begin{proof}
(a) By Theorem~\ref{thm:op} it suffices to identify the single-letter
union. Fix an admissible channel and let $D_j=\E(Y_j-\hY_j)^2$. Each
induced channel $p(\hat y_j\mid t_j)$ satisfies $\hY_j-T_j-S_j$, because
given $T_j$ the other coordinates and the channel randomness are
independent of $U_j$. The pairs $(T_j,S_j)$ are independent, so the chain
rule gives
\[
I(T;\hYv)=\sum_jI(T_j;\hYv,T^{j-1})\ge\sum_jI(T_j;\hY_j),
\]
\[
I(T;\hYv\mid S)=\sum_j\bigl[I(T_j;\hYv,S,T^{j-1})-I(T_j;S_j)\bigr]
\ge\sum_jI(T_j;\hY_j\mid S_j),
\]
using $I(T_j;S,T^{j-1})=I(T_j;S_j)<\infty$. If some $D_j=0$ the pair is
infinite and there is nothing to prove. Otherwise each pair
$\bigl(I(T_j;\hY_j),I(T_j;\hY_j\mid S_j)\bigr)$ lies in $\RW_j(D_j)$, and
the regions are upward closed. Conversely, a product of coordinate
channels attains the sum of their pairs with equality, by independence,
and each $\RW_j(D_j)$ is the union of quadrants above its single-letter
pairs \cite[Theorem~14(b)]{bond2026tit}.

(b) By (a) and \cite[Theorem~14(a)]{bond2026tit},
$L(D)=\min\sum_jL_j(D_j)$. Each $L_j$ is convex and nonincreasing on
$[0,\infty)$, strictly convex on $(0,\sigma_j^2)$ by Lemma~\ref{lem:slope},
with left derivative $-\kappa_j$ at $\sigma_j^2$ and zero beyond. For a
separable convex program with one linear constraint, an allocation is
optimal if and only if some $\theta\ge0$ lies in $-\partial L_j(D_j)$ for
every $j$, with $\theta(\sum_jD_j-D)=0$. Since $D<\sum_j\sigma_j^2$, $\theta=0$ is impossible, so
$\theta>0$, the constraint is active, and $D_j=D_j(\theta)$; no $D_j$ is
$0$ because $-L_j'\to\infty$ there. Each $D_j(\theta)$ is continuous and
nonincreasing, and strictly decreasing for $\theta>\kappa_j$, so
$\sum_jD_j(\theta)$ falls continuously and strictly from $\sum_j\sigma_j^2$
at $\theta=\min_j\kappa_j$ to $0$ as $\theta\to\infty$, and $\theta^\star$
is unique.

(c) $D_j(\theta)<\sigma_j^2$ exactly when $\theta>\kappa_j$, and
$\theta^\star$ increases as $D$ decreases, so coordinates enter in
increasing order of $\kappa_j$. For the rate the same argument with
$\tfrac12\log^{+}(\sigma_j^2/D_j)$ in place of $L_j$ gives levels
$1/(2\ln2\,\sigma_j^2)$ and reverse water-filling.
\end{proof}

\begin{corollary}\label{cor:parallelmisalign}
Let $m=2$ with $\sigma_1^2>\sigma_2^2$ and
$\sigma_1^2(1+\rho_1^2/\tau_1^2)<\sigma_2^2(1+\rho_2^2/\tau_2^2)$. For
every $D$ in the nonempty interval
$\bigl[\max\{2\sigma_2^2,\ \sigma_1^2+D_2(\kappa_1)\},\ \sigma_1^2+\sigma_2^2\bigr)$
the rate-optimal description describes coordinate $1$ alone and the
content-optimal description describes coordinate $2$ alone.
\end{corollary}

\begin{proof}
Reverse water-filling leaves coordinate $2$ undescribed and describes
coordinate $1$ exactly when $\sigma_2^2\le\lambda<\sigma_1^2$, that is,
for $2\sigma_2^2\le D<\sigma_1^2+\sigma_2^2$. The hypothesis is
$\kappa_2<\kappa_1$, and Theorem~\ref{thm:parallel}(b) describes
coordinate $2$ alone exactly when $\kappa_2<\theta^\star\le\kappa_1$, that
is, for $\sigma_1^2+D_2(\kappa_1)\le D<\sigma_1^2+\sigma_2^2$. Both
intervals end at $\sigma_1^2+\sigma_2^2$ and begin strictly below it.
\end{proof}

An instance makes the size concrete. Take
$(\sigma_1^2,\rho_1^2,\tau_1^2)=(1,0,1)$, a coordinate whose context is
useless, and $(\sigma_2^2,\rho_2^2,\tau_2^2)=(0.6,0.9,0.05)$, a smaller
coordinate with an informative context read in little noise. Then
$\kappa_1=0.7213$ and $\kappa_2=0.0633$, and the two descriptions describe
disjoint coordinates for every $D\in[1.2002,1.6)$. At $D=1.4$ the
rate-optimal description pays $R=L=0.1610$ bits, all on coordinate~$1$,
while the content-optimal description pays $L=0.0197$ bits and
$R=0.3272$ bits, all on coordinate~$2$: the first pays about eight times
the minimal content, the second about twice the minimal rate.

The theorem applies whenever an orthogonal change of coordinates of $Y$
and invertible linear changes of coordinates of $V$ and of $S$ reduce the
source to independent triples of this form, with the transformed side
information again of the form $S_j=V_j+U_j$ with independent $U_j$, since squared error is
invariant under an orthogonal map applied to $Y$ and $\hYv$ together and
both information coordinates are invariant under invertible maps of each
argument.

\section{Coupled Described Variables}
\label{sec:coupled}

Return to the general model of Section~\ref{sec:model}. By
Proposition~\ref{prop:posterior} the conditional content depends on a
description only through $\Sigma_e$, the uncertainty about $T$ that
remains for the party holding $S$ after it reads the description, and in
that coordinate the content endpoint is a concave maximization. Write
$h(\Sigma_e)=(\Sigma_e^{-1}-J_S)^{-1}$, which is $\Sigma_{e0}$.

\begin{theorem}\label{thm:maxdet}
For $0<D<\tr\Sigma_Y$,
\begin{equation}
L(D)=\tfrac12\log\det\Sigma_{T\mid S}
-\tfrac12\max\bigl\{\log\det\Sigma_e:\
0\prec\Sigma_e\preceq\Sigma_{T\mid S},\
\tr\bigl(Fh(\Sigma_e)F^{\mathsf T}\bigr)\le D\bigr\},
\label{eq:maxdet}
\end{equation}
and $h$ is matrix-convex on $\{0\prec\Sigma_e\preceq\Sigma_{T\mid S}\}$.
For a symmetric $Z$, $Fh(\Sigma_e)F^{\mathsf T}\preceq Z$ if and only if
\begin{equation}
\begin{bmatrix}
Z-F\Sigma_eF^{\mathsf T}&F\Sigma_eJ_S^{1/2}\\
J_S^{1/2}\Sigma_eF^{\mathsf T}&I-J_S^{1/2}\Sigma_eJ_S^{1/2}
\end{bmatrix}\succeq0,
\label{eq:lmi}
\end{equation}
so \eqref{eq:maxdet} is a determinant maximization under linear matrix
inequalities \cite{vandenberghe1998}. The maximizer $\Sigma_e^\star$ is
unique and the distortion constraint is active there. A feasible
$\Sigma_e$ is the maximizer if and only if
$\tr\bigl(Fh(\Sigma_e)F^{\mathsf T}\bigr)=D$ and there are $\varpi\ge0$ and
$\Theta\succeq0$ with
\begin{equation}
\tfrac12\Sigma_e-\varpi\,\Sigma_{e0}F^{\mathsf T}F\Sigma_{e0}
=\Sigma_e\Theta\Sigma_e,
\qquad
\Theta\bigl(\Sigma_{T\mid S}-\Sigma_e\bigr)=0,
\qquad\Sigma_{e0}=h(\Sigma_e),
\label{eq:kktcoupled}
\end{equation}
and then $\varpi>0$.
\end{theorem}

\begin{proof}
\emph{Change of variables.} The map $\Sigma_{e0}\mapsto\Sigma_e$ carries
$\{0\prec\Sigma_{e0}\preceq\Sigma_T\}$ onto
$\{0\prec\Sigma_e\preceq\Sigma_{T\mid S}\}$: $\Sigma_{e0}\preceq\Sigma_T$ if
and only if $\Sigma_{e0}^{-1}+J_S\succeq\Sigma_T^{-1}+J_S$, and conversely
$\Sigma_e\preceq\Sigma_{T\mid S}$ gives
$\Sigma_e^{-1}-J_S\succeq\Sigma_T^{-1}\succ0$. With
Proposition~\ref{prop:posterior} and Theorem~\ref{thm:op} this gives
\eqref{eq:maxdet}.

\emph{The matrix inequality and convexity.} The matrix
$J_S^{1/2}\Sigma_{T\mid S}J_S^{1/2}$ has the same nonzero eigenvalues as

\[
\Sigma_{T\mid S}^{1/2}J_S\Sigma_{T\mid S}^{1/2}\prec
\Sigma_{T\mid S}^{1/2}(\Sigma_T^{-1}+J_S)\Sigma_{T\mid S}^{1/2}=I,
\]
so on
the feasible set $J_S^{1/2}\Sigma_eJ_S^{1/2}\prec I$, and the Woodbury
identity gives
\[
h(\Sigma_e)=\Sigma_e+\Sigma_eJ_S^{1/2}\bigl(I-J_S^{1/2}\Sigma_eJ_S^{1/2}\bigr)^{-1}
J_S^{1/2}\Sigma_e .
\]
By the Schur-complement criterion \cite[Sec.~A.5.5]{boyd2004},
$Fh(\Sigma_e)F^{\mathsf T}\preceq Z$ if and only if \eqref{eq:lmi} holds.
The same criterion with $F$ replaced by $I_p$ and $Z$ a symmetric
$p\times p$ matrix shows that $h(\Sigma_e)\preceq Z$ if and only if a block
matrix affine in $(\Sigma_e,Z)$ is positive semidefinite. Averaging two
such instances, with $Z_i=h(\Sigma_{e,i})$, gives
$h\bigl((1-t)\Sigma_{e,0}+t\Sigma_{e,1}\bigr)\preceq(1-t)h(\Sigma_{e,0})+th(\Sigma_{e,1})$.

\emph{Existence, uniqueness, activity.} For small $\epsilon>0$ the
description with $\Sigma_{e0}=\epsilon\Sigma_T$ has distortion
$\epsilon\tr\Sigma_Y<D$, so the feasible set has an interior point. The
objective is bounded above on the bounded set
$\Sigma_e\preceq\Sigma_{T\mid S}$, its superlevel sets are closed and
bounded away from singular matrices, and so the maximum is attained;
$\log\det$ is strictly concave and the feasible set convex, so the
maximizer is unique. If the distortion were slack at the maximizer,
moving toward $\Sigma_{T\mid S}$ would stay feasible and strictly raise
$\log\det\Sigma_e$: $\Sigma_{T\mid S}$, with distortion $\tr\Sigma_Y>D$, is
infeasible, so $\Sigma_e\ne\Sigma_{T\mid S}$ and
$\tr\bigl(\Sigma_e^{-1}(\Sigma_{T\mid S}-\Sigma_e)\bigr)>0$.

\emph{Optimality conditions.} Slater's condition holds, so the
Karush--Kuhn--Tucker conditions, stationarity, complementary slackness for
both constraints, and nonnegativity of the multipliers, are necessary and
sufficient. The differential of $h$ is
$dh=h\,\Sigma_e^{-1}\,d\Sigma_e\,\Sigma_e^{-1}h$, so the gradient of
$\tr(FhF^{\mathsf T})$ is
$\Sigma_e^{-1}\Sigma_{e0}F^{\mathsf T}F\Sigma_{e0}\Sigma_e^{-1}$, and
stationarity of
$-\tfrac12\ln\det\Sigma_e+\varpi\,\tr(FhF^{\mathsf T})
+\tr\bigl(\Theta(\Sigma_e-\Sigma_{T\mid S})\bigr)$, multiplied on both
sides by $\Sigma_e$, is the first condition of \eqref{eq:kktcoupled}
(the factor $\ln2$ is absorbed in $\varpi$). When the distortion constraint
is active, its complementary slackness holds for every $\varpi\ge0$, which
gives the stated test; the maximizer has an active constraint, so the test
is also necessary. If $\varpi=0$, the first condition gives
$\Theta=\tfrac12\Sigma_e^{-1}\succ0$, and the second forces
$\Sigma_e=\Sigma_{T\mid S}$, which is infeasible; so $\varpi>0$.
\end{proof}

When the optimal description has one dimension, the program has a
closed form. Write $\Delta=\tr\Sigma_Y-D$ for the reduction in error the
description must buy, and
\[
A(\Delta)=\Sigma_TF^{\mathsf T}F\Sigma_T-\Delta\,\Sigma_T .
\]
A one-dimensional description is $W=b^{\mathsf T}T+N$ with $b\ne0$ and
$N\sim\mathcal N(0,\sigma^2)$ independent, reproduced by $\E[Y\mid W]$.
Its read direction is $b$, and the direction in which it reduces the
error of the reproduction is $F\Sigma_Tb$.

\begin{theorem}\label{thm:onedim}
(a) A one-dimensional description meets the distortion constraint with
finite content if and only if $\Delta<\lambda_{\max}(\Sigma_Y)$, and the
least content among such descriptions is
\begin{equation}
L_{\mathrm{1d}}(D)=\tfrac12\log\Bigl(1+\frac{\Delta}{\mu_{\max}(\Delta)}\Bigr),
\label{eq:onedim}
\end{equation}
where $\mu_{\max}(\Delta)>0$ is the largest root of
$\det\bigl(A(\Delta)-\mu\,\Sigma_{T\mid S}\bigr)=0$. It is attained with $b$
a corresponding generalized eigenvector and
$\sigma^2=|F\Sigma_Tb|^2/\Delta-b^{\mathsf T}\Sigma_Tb$, at which the
distortion constraint is active.

(b) Suppose $\Delta<\lambda_{\max}(\Sigma_Y)$, and let $\Sigma_{e0}$ and $\Sigma_e$ be the error covariances of the
description of (a), and
$\Gamma(\varpi)=\tfrac12\Sigma_e-\varpi\,\Sigma_{e0}F^{\mathsf T}F\Sigma_{e0}$.
Then $L(D)=L_{\mathrm{1d}}(D)$ if and only if there is $\varpi\ge0$ with
$\Gamma(\varpi)\succeq0$ and $\Gamma(\varpi)b=0$, and then
$\varpi=\tfrac12b^{\mathsf T}\Sigma_eb/|F\Sigma_{e0}b|^2$.

(c) For $m=1$, $L(D)=L_{\mathrm{1d}}(D)$ at every $0<D<\operatorname{Var}Y$
\cite[Proposition~24]{bond2026tit}. In the normalization of
\cite[Section~III]{bond2026tit}, with scalar context,
$\operatorname{Var}Y=\operatorname{Var}V=1$, and $\Sigma_U=\tau^2$, the
closed form of \cite[Theorem~14(a)]{bond2026tit} is
$g^\star=1+(1-D)/\mu_{\max}(1-D)$.
\end{theorem}

\begin{proof}
(a) For the description $W$,
$\Cov(T\mid W)=\Sigma_T-\Sigma_Tbb^{\mathsf T}\Sigma_T/(b^{\mathsf T}\Sigma_Tb+\sigma^2)$,
so the distortion of $\E[Y\mid W]$ is
$\tr\Sigma_Y-|F\Sigma_Tb|^2/(b^{\mathsf T}\Sigma_Tb+\sigma^2)$, and the
constraint reads $\sigma^2\le|F\Sigma_Tb|^2/\Delta-b^{\mathsf T}\Sigma_Tb$.
The content
$I(T;W\mid S)=\tfrac12\log\bigl(1+b^{\mathsf T}\Sigma_{T\mid S}b/\sigma^2\bigr)$
decreases in $\sigma^2$, so at the largest admissible $\sigma^2$, where the
constraint is active, it is
$\tfrac12\log\bigl(1+\Delta\,b^{\mathsf T}\Sigma_{T\mid S}b/b^{\mathsf T}A(\Delta)b\bigr)$,
finite exactly when $b^{\mathsf T}A(\Delta)b>0$. Minimizing over $b$ is
maximizing the Rayleigh quotient
$b^{\mathsf T}A(\Delta)b/b^{\mathsf T}\Sigma_{T\mid S}b$ of a symmetric
pencil with $\Sigma_{T\mid S}\succ0$, whose maximum is $\mu_{\max}$. Some
$b$ has $b^{\mathsf T}A(\Delta)b>0$ if and only if
$\Sigma_T^{1/2}F^{\mathsf T}F\Sigma_T^{1/2}\not\preceq\Delta I$, and the
nonzero eigenvalues of that matrix are those of
$F\Sigma_TF^{\mathsf T}=\Sigma_Y$.

(b) The description of (a) has an active distortion constraint and
$\Sigma_e^{-1}=\Sigma_{T\mid S}^{-1}+bb^{\mathsf T}/\sigma^2$, so
$\Sigma_{T\mid S}-\Sigma_e$ is a positive multiple of
$\Sigma_{T\mid S}bb^{\mathsf T}\Sigma_{T\mid S}$ and
$\Sigma_e^{-1}\Sigma_{T\mid S}b$ is a positive multiple of $b$. With
$\Theta=\Sigma_e^{-1}\Gamma(\varpi)\Sigma_e^{-1}$, the conditions
\eqref{eq:kktcoupled} at this point therefore read $\Gamma(\varpi)\succeq0$
and $\Gamma(\varpi)b=0$, and Theorem~\ref{thm:maxdet} makes them necessary
and sufficient. Taking $b^{\mathsf T}\Gamma(\varpi)b=0$ gives $\varpi$, and
$F\Sigma_{e0}b\ne0$ because $\Gamma(\varpi)b=0$ with $\Sigma_e\succ0$.

(c) The first statement is \cite[Proposition~24]{bond2026tit}. For the
second, substituting $\mu=(1-D)/(g-1)$ into
$\det(A(1-D)-\mu\Sigma_{T\mid S})$ gives
$-(1-D)(1-\rho^2)P(g)/\bigl(s(g-1)^2\bigr)$, a computation confirmed by
computer algebra (Appendix~\ref{app:numver}). The root $g^\star>1$ of $P$
corresponds to the unique positive $\mu$, which is $\mu_{\max}$; the other
root of $P$ lies in $(0,1)$ and corresponds to a negative $\mu$.
\end{proof}

For the rate, the same computation with $J_S=0$ makes
$\Sigma_{T\mid S}=\Sigma_T$ and the read direction the principal axis of
$\Sigma_Y$, whatever the budget (the rate-optimal description itself is
one-dimensional when $\Delta\le\lambda_1(\Sigma_Y)-\lambda_2(\Sigma_Y)$). For the content the budget enters the
pencil.

\begin{theorem}\label{thm:turning}
Let $\mu_{\max}(\Delta)$ be simple and positive on an open interval
$\mathcal I\subseteq(0,\lambda_{\max}(\Sigma_Y))$, with generalized eigenvector $b(\Delta)$. Then $b(\Delta)$
spans the same line for every $\Delta\in\mathcal I$ if and only if
$b(\Delta_0)=F^{\mathsf T}\beta$ for some $\Delta_0\in\mathcal I$, where
$\beta$ is an eigenvector of $\Sigma_Y$ with $\Cov(V,\beta^{\mathsf T}Y)=0$.
Otherwise $b(\Delta)$ is constant on no open subinterval of $\mathcal I$.
\end{theorem}

\begin{proof}
A simple eigenvalue of a symmetric pencil with positive definite second
matrix, and its eigenvector, are analytic in the parameter
\cite{kato1995}. Write $A(\Delta)=A_0-\Delta\Sigma_T$.

Suppose $b$ is constant on an open subinterval. Differentiating
$(A_0-\Delta\Sigma_T)b=\mu_{\max}(\Delta)\Sigma_{T\mid S}b$ gives
$\Sigma_Tb=c\,\Sigma_{T\mid S}b$ for a scalar $c$, and
$c=b^{\mathsf T}\Sigma_Tb/b^{\mathsf T}\Sigma_{T\mid S}b>0$. Then
$A_0b=(\Delta+\mu_{\max}/c)\Sigma_Tb$, that is,
$F^{\mathsf T}F\Sigma_Tb=\lambda b$ with $\lambda=\Delta+\mu_{\max}/c>0$.
So $b=F^{\mathsf T}\beta$ with $\beta=F\Sigma_Tb/\lambda$, and
$\Sigma_Y\beta=F\Sigma_TF^{\mathsf T}\beta=\lambda\beta$. Multiplying
$\Sigma_Tb=c\,\Sigma_{T\mid S}b$ by
$\Sigma_{T\mid S}^{-1}=\Sigma_T^{-1}+J_S$ gives $J_S\Sigma_Tb=(c-1)b$. The
left side lies in the range of $H^{\mathsf T}$ and the right side in the
range of $F^{\mathsf T}$, so both vanish:
$\Sigma_U^{-1}\Cov(V,Y)\beta=0$, that is, $\Cov(V,\beta^{\mathsf T}Y)=0$.
This proves the last sentence of the theorem and the forward direction.

Conversely, let $b_0=F^{\mathsf T}\beta$ with $\Sigma_Y\beta=\lambda\beta$
and $\Cov(V,\beta^{\mathsf T}Y)=0$. Then $\Sigma_Tb_0=\lambda b_0$ and
$J_S\Sigma_Tb_0=0$, so $\Sigma_{T\mid S}b_0=\Sigma_Tb_0=\lambda b_0$ and
$A(\Delta)b_0=(\lambda-\Delta)\Sigma_{T\mid S}b_0$ for every $\Delta$: $b_0$
is a generalized eigenvector with eigenvalue $\lambda-\Delta$ throughout.
Let $\mathcal E=\{\Delta\in\mathcal I:\lambda-\Delta=\mu_{\max}(\Delta)\}$.
It contains $\Delta_0$ and is closed in $\mathcal I$ by continuity. It is
open: at $\Delta_1\in\mathcal E$ the eigenvalue $\mu_{\max}(\Delta_1)$ is
simple, so for $\Delta$ near $\Delta_1$ exactly one eigenvalue of the
pencil lies near it, namely $\mu_{\max}(\Delta)$, and $\lambda-\Delta$ is an
eigenvalue near $\lambda-\Delta_1=\mu_{\max}(\Delta_1)$. Since
$\mathcal I$ is connected, $\mathcal E=\mathcal I$, and by simplicity
$b(\Delta)$ spans the line of $b_0$ on all of $\mathcal I$.
\end{proof}

An instance shows the effect. Take $\Sigma_Y=\mathrm{diag}(1,0.4)$, a
scalar context with $\operatorname{Var}V=1$ and $\Cov(Y,V)=0.6\,u$ for the
unit vector $u$ at $50^\circ$, and $\Sigma_U=\tau^2=0.05$. The certificate
of Theorem~\ref{thm:onedim}(b) holds for every $D\ge D_\star$, with
$D_\star=1.08530$ to five decimal places, out of $\tr\Sigma_Y=1.4$. The
rate-optimal description reduces the error along the first axis at every
$D\ge0.8$, in particular throughout the one-dimensional regime. The
content-optimal description reduces it along a direction at
$47.48^\circ$ at $D=1.38$, $46.76^\circ$ at $D=1.3$, $45.07^\circ$ at
$D=1.2$, and $41.15^\circ$ at $D=1.1$. The direction $F\Sigma_Tb$ is a function of
the read direction $b$, so the read direction turns as well.

The direction is not determined by the law of $(Y,S)$ either, which is
the analogue for directions of \cite[Corollary~21]{bond2026tit}. Four
sources with this $\Sigma_Y$, the context along the same $u$, and
$|\Cov(Y,V)|^2/(1+\tau^2)=0.3$ share the law of $(Y,S)$ up to scaling of
$S$. At $(|\Cov(Y,V)|^2,\tau^2)=(0.33,0.1)$, $(0.36,0.2)$, $(0.39,0.3)$,
and $(0.42,0.4)$, at $D=1.3$ and inside the one-dimensional regime of
each, their content-optimal descriptions reduce the error along
$42.83^\circ$, $37.61^\circ$, $33.75^\circ$, and $30.84^\circ$, at
contents $0.02366$, $0.03525$, $0.04206$, and $0.04649$ bits.

Below $D_\star$ the optimal description has two dimensions, and its value
is the solution of \eqref{eq:maxdet}, which interior-point methods compute
to any accuracy.

\section{Discussion}
\label{sec:discussion}

For a scalar described variable the content-optimal description differs
from the rate-optimal one in how much it loads on the context
\cite{bond2026tit}. For a described vector it differs in which parts of
the vector it resolves. With independent coordinates it can resolve
different coordinates altogether, and with coupled coordinates it reads a
direction that moves with the budget, set by a pencil that holds the
decoder's distortion, the holder's information, and the budget together.
The rate-optimal description, like principal components and the Gaussian
information bottleneck, reads fixed directions.

Several problems are open. A closed form for the coupled function below
the one-dimensional regime, and for the regime boundary $D_\star$ as a
function of the source, is not known. Both coordinates are convex in the error covariance $\Sigma_{e0}$:
$-\log\det\Sigma_{e0}$ is convex, and $\Sigma_e=(\Sigma_{e0}^{-1}+J_S)^{-1}$ is
matrix-concave in $\Sigma_{e0}$, so $-\log\det\Sigma_e$ is convex as well. Every
weighted combination of rate and content is therefore a convex function of
$\Sigma_{e0}$ under the linear constraint $\tr(F\Sigma_{e0}F^{\mathsf T})\le D$,
and the coupled rate--content frontier is computable by convex
programming; a closed-form description of its weighted minimizers is not
known.
Stationary vector sources, with spectral versions of the function, and
the finite-blocklength behavior of the pair $(R_n,L_n)$ are not treated.

\appendices
\section{Numerical Checks}
\label{app:numver}

The identities and instances were checked independently of the
derivations, on a remote workstation, by scripts that accompany the
manuscript sources. \texttt{verify\_parallel.py} confirms the slope and
curvature identities of Lemma~\ref{lem:slope} symbolically and against
finite differences, matches the equal-slope allocation of
Theorem~\ref{thm:parallel} against a brute-force search, checks that no
jointly Gaussian channel correlated across coordinates improves on it,
and computes the numbers of the parallel instance.
\texttt{verify\_coupled.py} confirms the scalar identity of
Theorem~\ref{thm:onedim}(c) symbolically, compares the program
\eqref{eq:maxdet} with a direct search over Gaussian descriptions on
random coupled sources, tests the certificate of
Theorem~\ref{thm:onedim}(b) against the optimizer, and computes the
numbers of the coupled instance, including $D_\star$ by bisection on the
certificate. The checks guard against algebraic slips; no proof depends
on them.

\section*{Acknowledgment}

In accordance with the IEEE policy on AI-generated text, the author
discloses the use of Anthropic's Claude models (Claude Opus 5.5), run
through the Claude Code environment in October 2026. The models drafted
the prose and the proofs from the author's statements of the problems,
carried out the numerical exploration and the prior-art search, wrote the
check scripts, and ran an independent referee pass on the proofs whose
findings were incorporated. The author is responsible for the content.

\begin{thebibliography}{99}
\bibitem{bond2026tit} A. H. Bond, ``Tradeoffs between rate and conditional
content with encoder-observed context,'' manuscript, 2026.
\bibitem{berger1971} T. Berger, \emph{Rate Distortion Theory: A Mathematical
Basis for Data Compression}. Englewood Cliffs, NJ: Prentice-Hall, 1971.
\bibitem{rahman2012} M. S. Rahman and A. B. Wagner, ``Rate region of the Gaussian
scalar-help-vector source-coding problem,'' \emph{IEEE Trans. Inf. Theory},
vol. 58, no. 1, pp. 172--188, 2012.
\bibitem{vandenberghe1998} L. Vandenberghe, S. Boyd, and S.-P. Wu, ``Determinant
maximization with linear matrix inequality constraints,'' \emph{SIAM J.
Matrix Anal. Appl.}, vol. 19, no. 2, pp. 499--533, 1998.
\bibitem{chechik2005} G. Chechik, A. Globerson, N. Tishby, and Y. Weiss,
``Information bottleneck for Gaussian variables,'' \emph{J. Mach. Learn.
Res.}, vol. 6, pp. 165--188, 2005.
\bibitem{steinberg2009} Y. Steinberg, ``Coding and common reconstruction,''
\emph{IEEE Trans. Inf. Theory}, vol. 55, no. 11, pp. 4995--5010, 2009.
\bibitem{lapidoth2014} A. Lapidoth, A. Mal\"ar, and M. Wigger, ``Constrained
source-coding with side information,'' \emph{IEEE Trans. Inf. Theory},
vol. 60, no. 6, pp. 3218--3237, 2014.
\bibitem{luxu2021} J. Lu and Y. Xu, ``Vector Gaussian Heegard-Berger problem
with common reconstructions,'' in \emph{Proc. Int. Conf. Wireless Commun. Signal
Process. (WCSP)}, 2021, pp. 1--5.
\bibitem{posner1975} E. C. Posner, ``Random coding strategies for minimum
entropy,'' \emph{IEEE Trans. Inf. Theory}, vol. 21, no. 4, pp. 388--391,
1975.
\bibitem{boyd2004} S. Boyd and L. Vandenberghe, \emph{Convex Optimization}.
Cambridge, U.K.: Cambridge Univ. Press, 2004.
\bibitem{kato1995} T. Kato, \emph{Perturbation Theory for Linear Operators},
reprint of the 1980 ed. Berlin, Germany: Springer, 1995.
\end{thebibliography}

\end{document}
